\section{What pretraining supplies to an agent}
\label{sec:C07:objective}
Autoregressive pretraining minimizes predictive loss on an encoded corpus, not task regret in an interactive environment. For sequences $x$ drawn from a mixture $\mathcal D$, the objective is $\mathcal L(\theta)=-\mathbb E_{x\sim\mathcal D}\sum_i w_i\log p_\theta(x_i\mid x_{<i})$, where fixed nonnegative weights $w_i$ specify which positions contribute. The model can acquire linguistic, procedural, and factual regularities useful to agents, but the loss does not distinguish an executed operation from a story about execution unless the data and interface preserve that distinction. A corpus containing fluent success narratives can teach the form of an answer without teaching reliable state verification.

Mixture design is therefore a policy-relevant intervention. Code, documentation, mathematical derivations, conversations, tool traces, and multimodal data contribute different conditional structures. Their sampling weights affect both information exposure and the optimization trajectory. A small collection of repeated tool traces can receive more effective training weight than a large rare domain. Counting unique documents is insufficient; report deduplication units, sampling probabilities, token exposure, and epoch repetition. For a stateful agent, training on histories with explicit failures, corrections, and permission boundaries may be more relevant than adding another large batch of generic prose.

Loss masks define the learning target in serialized interactions. A trace can contain user inputs, assistant actions, tool observations, and evaluator annotations. Training every position asks the model to predict all of them; training only assistant positions asks it to imitate the agent's behavior given the surrounding records. Neither is automatically correct for every objective. The decisive requirement is causal visibility. Future tool responses or retrospective summaries must not appear in the conditioning context for an earlier action. Otherwise the learned policy is rewarded for using evidence unavailable at execution time.

Sequence packing adds a related implementation hazard. Several documents can share one physical tensor for utilization, while an attention mask prevents unwanted cross-document conditioning. Omitting the mask changes the data distribution by letting one document predict another. For independent text this may merely alter modeling assumptions; for agent traces it can leak answers, test outcomes, or private task metadata across examples. Position resets, end-of-document markers, loss masking, and attention masking should be checked together. A shape-correct batch is not necessarily a semantically correct training example.

